
\begin{proposition}[The ruled obstruction with a reduced conductor image]\label{cdp:st:A.5.9}
Let $S\subset Y$ be an irreducible nonnormal reduced Cartier
divisor in a compact smooth complex threefold with
$H^2(Y,\Z)=0$ and $\chi(\OO_S)=0$.
Let $\eta:T\to S$ be its normalization and
$\pi:Z\to T$ its minimal resolution, and put $h=\eta\pi$.
Suppose that $Z$ is irrational ruled, with ruling $p:Z\to B$,
that every exceptional curve of $\pi$ is vertical, and that
the effective divisor of \ref{cdp:st:A.5.3} has the form
\[
 K_Z+D_Z=h^*\omega_S,\qquad D_Z=2C+V,                 \tag{A.RC.1}\label{cdp:eq:A.RC.1}
\]
where $C$ is a section in the strict conductor and $V$ is
vertical. Suppose also that the conductor image scheme
$\Gamma=\operatorname{Spec}_S(\OO_S/\mathfrak c)$ is reduced,
and that $H^0(S,K\otimes L|_S)\ne0$ for every $L\in\Pic(Y)$,
with $K$ as in \ref{cdp:st:A.1.1}.
Then every ambient line bundle pulls back trivially to $T$.
In particular
\[
 H^2\bigl(T,\omega_T\otimes
       \eta^*((K_Y^{-1}\otimes L)|_S)\bigr)\simeq\C
       \quad\hbox{for every }L\in\Pic(Y).            \tag{A.RC.2}\label{cdp:eq:A.RC.2}
\]
The proof forces $B$ and $\Gamma$ to be smooth elliptic curves,
$Z=T$ to be minimal ruled, and $D_Z=2C$ with $C^2=0$.
Thus no local-complete-intersection assumption on the original
conductor image curve is needed in this reduced-scheme case.
\end{proposition}
\begin{proof}
Write $\mathcal B=\eta_*\OO_T$ and
$\mathcal Q=\mathcal B/\OO_S$. We first specify conductor
duality, including its twists. Since $S$ is Gorenstein of
dimension two and $\mathcal B$ is maximal Cohen--Macaulay,
its dual is $\mathfrak c=\mathcal Hom_S(\mathcal B,\OO_S)$,
and maximal Cohen--Macaulay biduality gives
$\mathcal Hom_S(\mathfrak c,\OO_S)=\mathcal B$.
The quotient $\OO_\Gamma$ is Cohen--Macaulay of dimension one:
this follows from the depth lemma in
$0\to\mathfrak c\to\OO_S\to\OO_\Gamma\to0$.
Applying $\mathcal Hom_S(-,\OO_S)$ to that sequence, its
initial terms identify the actual normalization inclusion and
its cokernel as
\[
 0\longrightarrow\OO_S\longrightarrow\mathcal B
 \longrightarrow
 \mathcal Ext^1_S(\OO_\Gamma,\OO_S)\longrightarrow0,
 \qquad
 \mathcal Q\simeq\omega_\Gamma\otimes\omega_S^{-1}|_\Gamma.
                                                               \tag{A.RC.3}\label{cdp:eq:A.RC.3}
\]
Adjunction makes $\omega_S$ the restriction of
$K_Y\otimes\OO_Y(S)$. Its degree is zero on every compact
curve, because $H^2(Y,\Z)=0$. Curve Riemann--Roch and Serre
duality therefore give
$\chi(\mathcal Q)=\chi(\omega_\Gamma)=-\chi(\OO_\Gamma)$.
On the other hand, \ref{cdp:st:A.5.8} gives rational singularities of $T$
and $\chi(\OO_T)=1-g(B)$. The normalization sequence and
$\chi(\OO_S)=0$ now give the exact numerical identity
\[
                   \chi(\OO_\Gamma)=g(B)-1.          \tag{A.RC.4}\label{cdp:eq:A.RC.4}
\]

We next prove that $\Gamma$ is connected. Contract ruling
$(-1)$-curves from $Z$ to a minimal ruled bundle $Z_0\to B$.
Under a single contraction $\sigma:Z_1\to Z_2$ with exceptional
curve $E$, let $D_2=\sigma_*D_1$, let $a$ be the coefficient
of $E$ in $D_1$, and let $m$ be the multiplicity of $D_2$ at
the contraction point. The canonical divisor formula gives
\[
 K_{Z_1}+D_1=\sigma^*(K_{Z_2}+D_2)+(1+a-m)E.
                                                               \tag{A.RC.5}\label{cdp:eq:A.RC.5}
\]
Numerical triviality and intersection with $E$ force $a=m-1$.
The same numerical triviality descends to $Z_2$. On $Z_0$
the pushed divisor is $2C_0+V_0$, with $C_0$ a section and
$V_0$ a nonnegative sum of fibres, so its support is connected.
Reverse the contractions. A centre outside the existing support
would produce coefficient $-1$ on its exceptional curve; that
coefficient persists on strict transforms and contradicts the
effectiveness of the final $D_Z$. A centre of multiplicity one
is a smooth point of one coefficient-one component and its
exceptional coefficient is zero, so it does not disconnect the
support. At a centre of multiplicity at least two the exceptional
curve has positive coefficient and connects all strict-transform
branches through that point. Induction proves that
$\operatorname{Supp}D_Z$ is connected.

Every positive exceptional component of $D_Z$ maps into
$\Gamma$. Indeed away from $\Gamma$ the surface $T$ equals
the Gorenstein surface $S$ and has rational singularities by
\ref{cdp:st:A.5.8}. Such a singularity is crepant on its minimal resolution:
a generator of its dualizing sheaf pulls back holomorphically
by rational duality; its effective exceptional zero divisor
has nonnegative intersection with each exceptional curve by
adjunction and minimality, so negative definiteness makes it
zero, as in \ref{cdp:st:A.5.2}. Thus \eqref{cdp:eq:A.RC.1} has no positive exceptional
coefficient there. The strict conductor maps onto the whole
conductor support. It follows that
$h(\operatorname{Supp}D_Z)=\operatorname{Supp}\Gamma$,
which proves connectedness.

Let $D\subset\Gamma$ be the image of $C$. Localize \eqref{cdp:eq:A.RC.3}
at the generic point of $D$. Because $\Gamma$ is reduced,
$\OO_\Gamma$ there is the function field $k(D)$, and its
dualizing module has dimension one over that field. The
elementary conductor quotient sequence is
\[
 0\longrightarrow\OO_\Gamma\longrightarrow
          \mathcal B/\mathfrak c
   \longrightarrow\mathcal Q\longrightarrow0.
\]
It follows that $\mathcal B/\mathfrak c$ has dimension two
over $k(D)$. The conductor coefficient two on $C$ contributes
$2\deg(C\to D)$ to this dimension. Therefore this degree is
one and there is no other normalization branch over $D$.
The generic transverse ring has conductor $t^2\C\{t\}$ and
is exactly $\C+t^2\C\{t\}=\C\{t^2,t^3\}$, the cusp ring.
Thus $D^\nu\simeq C\simeq B$. This deduction also rules out
any additional vertical branch over the same image curve.

A connected compact reduced curve has only constant global
holomorphic functions. From \eqref{cdp:eq:A.RC.4} we therefore obtain
$h^1(\Gamma,\OO_\Gamma)=2-g(B)$.
Restriction to the reduced component $D$ gives a surjection
on $H^1$, because its kernel is a coherent sheaf on a curve
and has no $H^2$. Normalization of $D$ gives
$h^1(D,\OO_D)\ge g(B)$. Hence
\[
             2-g(B)\ge h^1(D,\OO_D)\ge g(B).         \tag{A.RC.6}\label{cdp:eq:A.RC.6}
\]
As $g(B)\ge1$, all quantities equal one. The normalization
quotient of $D$ has length zero, so $D$ itself is smooth
elliptic. Statement \ref{cdp:st:C2d} applies to this actual smooth curve
in $Y$ and makes the transversal cusp constant along it.
No other singular curve meets $D$. Connectedness of $\Gamma$
therefore gives $\Gamma=D$. Its normalization inverse image
is the single smooth curve $C$, the normalization is smooth
near it, and \ref{cdp:st:C2d} gives $C^2=0$.

Adjunction in \eqref{cdp:eq:A.RC.1} yields $V\cdot C=0$. A ruling
$(-1)$-curve outside $V$ cannot meet $C$, since its
intersection with $K_Z+2C+V$ would be at least $-1+2$.
A component of $V$ cannot meet $C$ either. Thus all
contractions to $Z_0$ avoid $C$ and preserve its square zero.
On the minimal bundle, $K_{Z_0}+2C_0$ is numerically trivial,
so its nonnegative vertical divisor $V_0$ is zero.
The first reverse blowup would be away from $C_0$ and have
coefficient $-1$ by \eqref{cdp:eq:A.RC.5}. There are no such blowups.
Thus $Z$ is minimal ruled and $D_Z=2C$.
An exceptional curve for $Z\to T$ would have negative square;
adjunction with $K_Z\equiv-2C$ would make it rational and
therefore vertical. A minimal ruled bundle has only fibres
of square zero, so no such exceptional curve exists and $Z=T$.

The constant-cusp computation \ref{cdp:st:A.2.4} now gives the injection
$K\hookrightarrow\eta_*(h^*\OO_Y(-S)\otimes\OO_Z(C))$.
Every ambient line is topologically trivial on $Z$ and is
the pullback of a degree-zero line on the elliptic curve $B$.
The section exact sequence, together with $R^1p_*\OO_Z=0$,
gives
\[
 0\longrightarrow\OO_B\longrightarrow p_*\OO_Z(C)
 \longrightarrow N_{C/Z}\longrightarrow0,
 \qquad\deg N_{C/Z}=0.                              \tag{A.RC.7}\label{cdp:eq:A.RC.7}
\]
After a degree-zero twist $P$, both end terms have zero
sections unless $P=\OO_B$ or $P=N_{C/Z}^{-1}$.
The image of $\Pic(Y)$ in $\Pic^0(B)$ is connected:
$H^2(Y,\Z)=0$ makes $\Pic(Y)$ the connected image of
$H^1(Y,\OO_Y)$ under the exponential map.
All-twist persistence of $K$, and translation by the ambient
line $\OO_Y(-S)$, put this entire image in those at most two
points. It is therefore the identity. This proves triviality
of every ambient pullback. Serre duality on $T$ proves
\eqref{cdp:eq:A.RC.2}.
\end{proof}


